%\subsection{Measure Theoretic Probability Theory}%

%
\subsection{Crash Course: Measure Theoretic Probability}%

\subsubsection{Elementary Probability Theory - Recap and Beyond}

\begin{Eg}[Winning a pie with a biased die]%
    \label{example-die}
    You are allowed to roll a biased die with 6 sides. If you roll a $5$ or $6$ you win a \emph{car}, a $4$ 
    gives you a \emph{mug} 
    and $1,2,3$ wins you an apple \emph{pie}. 
    In this case the sample space is $\Wcal :=\{1,2,3,4,5,6\}$ and the die introduces a probability distribution $P$ on $\Wcal$.
    Since the die is biased, we have to specify 
    %for the probability distribution $P$
    each of the probability masses to throw those numbers separately: 
    \begin{align*}
        p(1) & = 0.5, &
        p(2) & = p(3) = p(4) = 0.1, &
        p(5) & = 0.15 , & 
        p(6) & = 0.05.
    \end{align*}
    We are now interested in the probabilities of the events of winning those 3 different prices. 
    For this we consider the 'prize' space: $\Zcal :=\{\mathrm{pie}, \mathrm{mug}, \mathrm{car}\}$. 
    We can then formalize the outcome via the map $F$:
    \[\begin{array}{cccc}
        F: & \Wcal &\to& \Zcal,\\
           &1,2,3 &\mapsto& \mathrm{pie},\\
        & 4 & \mapsto& \mathrm{mug},\\
        & 5,6 & \mapsto& \mathrm{car}.
    \end{array}\]
    To compute the probability of winning each of the prizes we need to 'push' the probability distribution $P$, which lives on the space $\Wcal$, to the space $\Zcal$. We can do this as follows:
    \begin{align*}
        P(F=\mathrm{pie}) &= P(F^{-1}(\{\mathrm{pie}\})) && = P(\{1,2,3\}) && = p(1)+p(2)+p(3) && = 0.7,\\
            P(F=\mathrm{mug}) &= P(F^{-1}(\{\mathrm{mug}\})) && = P(\{4\}) && = p(4) && = 0.1,\\
            P(F=\mathrm{car}) &= P(F^{-1}(\{\mathrm{car}\})) && = P(\{5,6\}) && = p(5)+p(6) && = 0.2,
    \end{align*}
    where $F^{-1}(C) := \{w \in \Wcal\,|\,F(w) \in C\}$ is the \emph{pre-image} of $C \ins \Zcal$.
\end{Eg}

\begin{Disc}%
    The simple example \ref{example-die} already provides us with the main examples for the typical 
    probability-theoretic terminology and important insights:
    \begin{enumerate}
        \item We call the tuple $(\Wcal,P)$ a \emph{probability space}. It is important to note that $P$ was defined on $\Wcal$, not $\Zcal$.
        \item We call the map $F$ a \emph{random variable}, which is really nothing else than a map from a probability space to another space.
        \item \emph{Events} are modelled by \emph{subsets} $B \ins \Wcal$, not just by single elements $w \in \Wcal$. For example consider the event that you don't win a car. This event can't be represented by a single element in $\Wcal$ or $\Zcal$.
        \item In this example we can compute the probability of an event by \emph{additivity} of $P$ and the use of the probability mass function,
            via $P(B) = \sum_{ w \in B} p(w)$.
        \item The distribution of the prizes, i.e.\ the distribution of random variable $F$, assigns 
               probabilities to events $C \ins \Zcal$ and can be computed using the \emph{pre-image} of $F$ 
               via $P(F \in C) = P(F^{-1}(C))$, where the latter is now an event $F^{-1}(C) \ins \Wcal$, 
               which we already know how to deal with. 
           \item The distribution of $F$ on $\Zcal$ here is also called the \emph{push-forward} distribution or \emph{image} distribution 
            of $P$ via $F$ or just the \emph{law} of $F$. It is often abbreviated as: $P_F$, $P^F$, $F_*P$ or $P(F)$. 
            Again note: $P(F)(C) := P(F \in C)= P(F^{-1}(C))$.
        \item So $(\Zcal,P(F))$ forms a probability space on its own and as soon as we know $P(F)$ we don't need any information about $(\Wcal,P)$ anymore if all we are interested in is the events in $\Zcal$ and the law of $F$. All randomness on $\Zcal$ is fully specified by $P(F)$.
    \end{enumerate}
\end{Disc}

\begin{Eg}%
    \label{example-normal}
    Now consider the standard normal distribution $\Ncal(0,1)$ on $\R$, which is given by the \emph{probability density} function:
    \[ p(w) = \frac{1}{\sqrt{2 \pi}} \exp\lp - \frac{1}{2} \cdot w^2 \rp.   \]
    The probability of an event $A \ins \R$ is then given by:
    \[ P(A) = \int_A p(w) \, dw,\]
    in case $A$ can be integrated (i.e.\ if it is not a too pathological set).
    For instance, if $A= [a,b] \cup [c,d]$ with $a \le b <c\le d$ we get:
    \[P(A) = \int_a^b p(w)\,dw + \int_c^d p(w)\,dw.\]
    Note that, even though $p(w) >0$ for every $w \in \R$, we have:
\[  P(\{x\}) = 0.\]
Now consider random variable $F: \R \to \R$ with $F(w)=\sin(w)$. 
It is not immediately clear how to define the probability distribution of $F$ when only working with probability densities. It is even more difficult to derive the probability density for $F$ in this setting.
\end{Eg}%

\begin{Disc}%
    \begin{enumerate}
        \item The examples \ref{example-die} and \ref{example-normal} show that many probability distributions can be represented either by \emph{probability mass} functions (\emph{discrete} case), $w \mapsto p(w)$, or \emph{probability density} functions (\emph{absolute continuous} case), $w \mapsto p(w)$.
        \item Both cases have in common that one only needs a function that takes elements $w \in \Wcal$ as arguments, in contrast to subsets $A \ins \Wcal$. This is usually the reason why only the discrete and absolute continuous cases are taught in e.g.\ Machine Learning classes.
        \item Note that, in the discrete case with $K$ classes, one only needs to specify the $K$ values $p(1),\dots,p(K)$, in contrast to the $2^K$ values on subsets $P(A)$ for $A \in 2^\Wcal$ (the \emph{power set} of $\Wcal$ consisting of all subsets $A \ins \Wcal$), as those values can be derived from $p(w)$ values using additivity.
        \item We have problems defining probability distributions of random variables for absolute continuous distributions when we are only allowed to work with probability densities.
        \item Measure theory is the framework that directly works with subsets $A \ins \Wcal$, in contrast to elements $w \in \Wcal$.

    \end{enumerate}
\end{Disc}%


\subsubsection{Why Measure Theory?}%
%
\paragraph{Discrete and absolute continuous distributions are not general enough}%

\begin{Eg}[Simple example of a non-discrete non-absolute-continuous distribution]%
    Consider a uniformly distributed random variable on the interval $\Xcal:=[0,1]$, i.e.\ $X \sim \Ucal[0,1]$, which has probability density:
    \[ p(x) = \I_{[0,1]}(x).\]
    Consider an exact copy of $X$, which we call $Y:=X$, on $\Ycal :=[0,1]$. Now consider the joint distribution of $(X,Y)$ 
    on $\Xcal \times \Ycal = [0,1]^2$.
    Then only values on the diagonal $\Delta:=\{ (x,x) \,|\, x \in [0,1]\}$ can be realized by $(X,Y)$.
    This simple distribution on $[0,1]^2$ is not discrete (as it can attain uncountably many values), and it is also not 
    absolute continuous, since we have: $\int_\Delta \,dx\,dy = 0$, i.e.\ the (2-dimensional) area of the (1-dimensional) line is zero. 
    This implies that any density function $p$ would satisfy: $\int_\Delta p(x,y)\,dx\,dy =0$ as well. This is in contrast to the fact that a probability distribution should always be normalized:
    \[ 1 = P((X,Y) \in \Delta) = \int_\Delta p(x,y)\,dx\,dy.\]
    Note that we don't need a probability density to be able to assign probabilities to subsets $D \ins [0,1]^2$.
    We can just use  the push-forward map:
    \[(X,Y):\; [0,1] \to [0,1] \times [0,1], \quad x \mapsto (x,x).\]
    and compute:
    \[P((X,Y) \in D) = P(\{x \in [0,1]\,|\, (x,x) \in D\}),\]
    where $P$ on the right here denotes the uniform distribution on $[0,1]$.
\end{Eg}%

\begin{Not}[Unifying the notations to measure theoretic ones]%
    Let $X$ be a  random variable taking values in space $\Xcal$ and with probability distribution $P$.
    Let $F: \Xcal \to \R$ be a function. Then we will change the notations for expectation values as follows.
 \begin{enumerate}
     \item Let $X$ be a \emph{discrete} random variable with \emph{probability mass function} $p$. Then define:
 \begin{align*}
     \E[F(X)] &= \sum_{x \in \Xcal} F(x) \cdot p(x) \\
         &=: \int F(x)  \, P(dx)\\
         &=: \int F(x)  \, dP(x)\\
         &=: \int F  \, dP.
 \end{align*}
  We will consider sums to be special cases of measure integrals.
     \item Let $X$ be a \emph{absolute continuous} random variable with \emph{probability density function} $p$. Then define:
 \begin{align*}
     \E[F(X)] &= \int_\Xcal F(x) \cdot p(x)\,dx \\
         &=: \int F(x)  \, P(dx)\\
         &=: \int F(x)  \, dP(x)\\
         &=: \int F  \, dP.
 \end{align*}
 \end{enumerate}
 So both cases can be unified with the 3 commonly used notations:
 \[ \E[F(X)] = \int F\,dP = \int F(x)\,dP(x) = \int F(x)\,P(dx).  \]
 Note that in both cases we also can write: $P(A) = \int \I_A \, dP$.
\end{Not}

\begin{Exc}
    Show that the following relation holds:
    \[ \int F(x)\,P(dx) = \int z \,P^F(dz).\]
\end{Exc}

\paragraph{Defining probability distributions on all subsets is too general}%

\begin{Rem}\label{rem-power-set}
    When we want to work with a (probability) measure $\mu$ we at least want to require that it is \emph{countably additive}, i.e.\ that for pairwise \emph{disjoint} subsets $A_n \ins \Xcal$, $n \in \N$, we have:
    \[\mu\lp\bigcup_{n \in \N} A_n\rp = \sum_{n \in \N} \mu(A_n).\]
%    where $\bigdcup$ is the symbol for disjoint union.
    We will see below that if we do not restrict the subsets $A_n$ in some way we will encounter strange behaviour.
\end{Rem}

\begin{Thm}[Vitali, non-existence of Lebesgue measure on all subsets]%
   \label{thm-non-lebesgue}
    There does NOT exist a measure $\lambda$ on $[0,1]$ such that:
    \begin{enumerate}
        \item $\lambda$ can measure every $A \in 2^{[0,1]}$, and:
        \item $\lambda\lp[a,b] \rp = b-a$ for all $a\le b \in [0,1]$.
    \end{enumerate}
    In other words, there does NOT exist a uniform distribution on $[0,1]$ that can consistently assign values to all subsets.\\
    But: such a measure with property 2. exists on the set $\Bcal_{[0,1]}$ of all so called \emph{Borel subsets} (or even Lebesgue subsets) 
    of $[0,1]$.
    %It is called the \emph{Lebesgue measure} on $\R$ (or, more precisely, on $\Bcal_\R$). 
    Similar statements hold for higher dimensions and $\R^D$ and higher dimensional volumes.
\end{Thm}

\begin{Eg}[Vitali set]
    Consider the following equivalence relation on $[0,1]$:
    \[ r_1 \sim r_2 \qquad :\iff \qquad r_2-r_1 \in \Q.\]
    Let $[0,1]/\mathord{\sim}$ be the set of equivalence classes.
    By the axiom of choice there exists a representative system $V \ins [0,1]$ for $[0,1]/\mathord{\sim}$.
    This means that the map:
    \[ V \to [0,1]/\mathord{\sim}, \qquad v \mapsto [v],\]
    is bijective. For $q \in \Q$ consider the subset:
    \[V_q :=  V+q := \lC v + q \,|\, v \in V \rC  \ins \R. \]
    Let $\Qcal := [-1,1] \cap \Q$.  Note that $[0,1]$ and $V_q$ are uncountable while $\Q$ and $\Qcal$ are countably infinite. 
    We then have the inclusions:
    \[ [0,1] \ins  \bigcup_{ q \in \Qcal} V_q \ins [-1,2].\]
    The right inclusion is clear as:
     \[V + [-1,1] \ins [0,1]+[-1,1] \ins [-1,2].\]
    For the left inclusion let $ x\in [0,1]$. By construction there exists a $v \in V$ such that $v \sim x$. 
    So $x - v \in \Q$. Since $x,v \in [0,1]$ we also have that $x-v \in [-1,1]$. So $q := x-v \in [-1,1] \cap \Q =\Qcal$.
    This shows that $ x \in V_q$ for a $q \in \Qcal$. Thus both inclusions are shown.\\
    If we now assumed that $V$ would be Lebesgue-measurable then every $V_q$ would be as well as a translated version of $V$.
    We then would get that: $\lambda(V_q) = \lambda(V)$ for every $q\in \Q$. So we would get:
    \[ 1 = \lambda \lp[0,1]\rp \le \lambda\lp\bigcup_{ q \in \Qcal} V_q \rp   \le \lambda \lp[-1,2]\rp =3  ,\]
    which implies:
    \[ [1,3] \ni \lambda\lp\bigcup_{ q \in \Qcal} V_q \rp = \sum_{ q\in \Qcal} \lambda(V_q) = 
    \sum_{ q\in \Qcal} \lambda(V), \]
    which is contradictory. Indeed, $\lambda(V) =0$ can be ruled out as the sum would sum up to $0 \notin [1,3]$.
    But also $\lambda(V)>0$ can be ruled out as this would sum up to $\infty \notin [1,3]$.
    So the \emph{Vitali set} $V$ can not be Lebesgue-measurable.
\end{Eg}




\begin{Thm}[Banach-Tarski paradox]
    \label{thm-banach-tarski}
    The 3-dimensional unit ball $B_1(z)=\{ x \in \R^3\,|\, \|x-z\| \le 1\}$ centered at $z \in \R^3$ can be partitioned into a finite number of disjoint sets $A_1,\dots,A_K$ (e.g.\ $K=5$) such that each can then be rotated and translated in $\R^3$ such that they form TWO 3-dimensional unit balls $B_1(y_1)$ and $B_1(y_2)$.\\
    Note that the unit balls have well-defined volume (i.e.\ 3-dimensional Lebesgue measure) and translation and rotations are very well behaved and preserve volume, while the subsets $A_k$ are very pathological (i.e.\ non-Lebesgue-measurable).
\end{Thm}


\begin{figure}[ht]
    \centering
\includegraphics[scale=0.33]{Banach-Tarski-Paradox}
\caption[Banach-Tarski paradox]{Illustration of the Banach-Tarski paradox.\footnotemark}
\label{fig:banach-tarski}
\end{figure}
\footnotetext{\href{https://en.wikipedia.org/wiki/Banach\%E2\%80\%93Tarski\_paradox
}{https://en.wikipedia.org/wiki/Banach-Tarski\_paradox} }



\textbf{$\Longrightarrow$ Measure theory is the unifying `safe space' for probability theory!}



